\documentclass[border=10pt]{standalone}
\usepackage{tikz,ctex}
\usetikzlibrary{positioning}

\begin{document}

\begin{tikzpicture}[
    % 全局设置
    node distance = 1.6cm and 1.4cm,
    every node/.style = {
        draw,
        rounded corners,
        thick,
        align = center,
        font = \sffamily\small,
        minimum width = 2.0cm,
        minimum height = 0.9cm
    },
    % 样式定义
    centerStyle/.style = {
        fill = blue!45,
        text = white,
        font = \sffamily\bfseries\large,
        line width = 2pt,
        minimum size = 2.8cm,
        circle
    },
    branchTitle/.style = {
        font = \sffamily\bfseries,
        line width = 1.5pt
    },
    subNode/.style = {
        fill = white,
        font = \sffamily\small,
        minimum width = 2.2cm,
        minimum height = 0.8cm
    },
    leafNode/.style = {
        fill = white,
        font = \sffamily\footnotesize,
        minimum width = 1.9cm,
        minimum height = 0.7cm
    },
    arrowStyle/.style = {
        ->,
        thick,
        >=stealth,
        shorten >=2pt
    }
]

    % ================= 1. 中心节点 =================
    \node[centerStyle] (center) {9.6.模型平均 \\ Model Averaging};

    % ================= 2. 上方分支：集成方法 (红) =================
    \node[branchTitle, fill=red!15, above=of center] (ensemble) {集成方法 \\ Ensembles};

    \node[subNode, above=of ensemble, xshift=-6cm] (committee) {委员会 \\ Committees};
    \node[subNode, above=of ensemble, xshift=-2cm] (bagging) {Bagging \\ 自助聚合};
    \node[subNode, above=of ensemble, xshift=2cm] (bootstrap) {自助采样 \\ Bootstrap};
    \node[subNode, above=of ensemble, xshift=6cm] (boost) {Boosting \\ 提升};

    % ================= 3. 右侧分支：Dropout (蓝) =================
    \node[branchTitle, fill=orange!15, right=of center] (dropout) {Dropout \\ 随机失活};

    \node[subNode, right=of dropout, yshift=4cm] (mask) {掩码向量 \\ Mask $R_i$};
    \node[subNode, right=of dropout, yshift=2cm] (prob) {保留概率 $\rho$ \\ Keep Probability};
    \node[subNode, right=of dropout, yshift=0cm] (prune) {剪枝网络 \\ Pruned Networks};
    \node[subNode, right=of dropout, yshift=-2cm] (mc) {蒙特卡洛 \\ Monte Carlo};
    \node[subNode, right=of dropout, yshift=-4cm] (scale) {权重缩放 \\ Weight Scaling};

    % ================= 4. 下方分支：理论分析 (绿) =================
    \node[branchTitle, fill=green!15, below=of center] (theory) {理论分析 \\ Theory};

    \node[subNode, below=of theory, xshift=-6cm] (bias) {偏差--方差 \\ Bias--Variance};
    \node[subNode, below=of theory, xshift=-2cm] (uncorr) {误差不相关 \\ Uncorrelated Errors};
    \node[subNode, below=of theory, xshift=2cm] (bayes) {贝叶斯视角 \\ Bayesian View};
    \node[subNode, below=of theory, xshift=6cm] (coadapt) {共适应抑制 \\ Reduce Co-adaptation};

    % ================= 5. 左侧分支：实践要点 (紫) =================
    \node[branchTitle, fill=purple!15, left=of center] (practice) {实践要点 \\ Practical Points};

    \node[subNode, left=of practice, yshift=3cm] (cost) {计算开销 \\ Computational Cost};
    \node[subNode, left=of practice, yshift=1cm] (noisy) {噪声梯度 \\ Noisy Gradients};
    \node[subNode, left=of practice, yshift=-1cm] (test) {测试阶段 \\ Test Phase};
    \node[subNode, left=of practice, yshift=-3cm] (broad) {广泛应用 \\ Broad Applicability};

    % ================= 连线逻辑 =================
    % 中心 -> 四大分支标题
    \draw[arrowStyle, red] (center) -- (ensemble);
    \draw[arrowStyle, blue] (center) -- (dropout);
    \draw[arrowStyle, green!60!black] (center) -- (theory);
    \draw[arrowStyle, purple] (center) -- (practice);

    % 分支标题 -> 子节点 (上方)
    \draw[arrowStyle, red!60] (ensemble) -- (committee.south);
    \draw[arrowStyle, red!60] (ensemble) -- (bagging.south);
    \draw[arrowStyle, red!60] (ensemble) -- (bootstrap.south);
    \draw[arrowStyle, red!60] (ensemble) -- (boost.south);

    % 分支标题 -> 子节点 (右侧)
    \draw[arrowStyle, blue!60] (dropout) -- (mask.west);
    \draw[arrowStyle, blue!60] (dropout) -- (prob.west);
    \draw[arrowStyle, blue!60] (dropout) -- (prune.west);
    \draw[arrowStyle, blue!60] (dropout) -- (mc.west);
    \draw[arrowStyle, blue!60] (dropout) -- (scale.west);

    % 分支标题 -> 子节点 (下方)
    \draw[arrowStyle, green!60!black] (theory) -- (bias.north);
    \draw[arrowStyle, green!60!black] (theory) -- (uncorr.north);
    \draw[arrowStyle, green!60!black] (theory) -- (bayes.north);
    \draw[arrowStyle, green!60!black] (theory) -- (coadapt.north);

    % 分支标题 -> 子节点 (左侧)
    \draw[arrowStyle, purple!60] (practice) -- (cost.east);
    \draw[arrowStyle, purple!60] (practice) -- (noisy.east);
    \draw[arrowStyle, purple!60] (practice) -- (test.east);
    \draw[arrowStyle, purple!60] (practice) -- (broad.east);

\end{tikzpicture}

\end{document}